\section{Conclusion}
\label{sec:conclusion}

This project compared feature-based MLPs and sequence-based GRUs for inferring physical parameters from simulated three-dimensional projectile trajectories. The restricted initial-velocity task was solved extremely accurately by an MLP using engineered features, with test RMSEs of $\SI{0.00640}{\meter\per\second}$, $\SI{0.00677}{\meter\per\second}$, and $\SI{0.00304}{\meter\per\second}$ for $v_x$, $v_y$, and $v_z$, respectively.
This result is primarily attributable to the inclusion of a finite-difference estimate of the initial velocity and therefore demonstrates the value of task-specific feature engineering. For the more demanding joint prediction of $11$ launch and environmental parameters, the GRU models generally outperformed the all-parameter MLP, particularly for wind, spin direction, drag, and Magnus strength.

On complete, noise-free trajectories, explicitly providing velocity and acceleration channels produced the most accurate overall parameter recovery. For example, relative to the position-only full-trajectory GRU, the derivative GRU reduced the test RMSE of $u_x$ from $\SI{2.67}{\meter\per\second}$ to $\SI{1.36}{\meter\per\second}$ and that of $\beta_D$ from $\SI{0.00300}{\per\meter}$ to $\SI{0.00127}{\per\meter}$.
Random-crop augmentation was essential when only a prefix of the flight was available. The crop-trained derivative model retained useful performance at short observation fractions and, at $f=0.1$, reduced the $v_x$ RMSE from $\SI{6.62}{\meter\per\second}$ for the model trained only on complete trajectories to $\SI{2.16}{\meter\per\second}$. The auxiliary impact-point objective did not consistently improve recovery of the physical parameters.

The Gaussian-noise experiments exposed the main weakness of the present approach. Although numerical derivatives are highly informative for clean inputs, differentiating perturbed positions amplifies measurement noise. For the derivative GRU, adding only $\sigma=\SI{0.01}{\meter}$ of positional noise increased the $v_x$ RMSE from $\SI{1.03}{\meter\per\second}$ to $\SI{21.0}{\meter\per\second}$. Crop augmentation did not prevent this degradation because none of the models was trained with noisy observations. Future work should therefore include noise augmentation matched to the measurement process and investigate smoothing, regularized differentiation, or architectures that infer kinematic information directly from positions.\cite{chartrand2011numerical}

The conclusions are limited to the simulated parameter ranges and physical assumptions used here. In addition, the same test split was used for checkpoint selection and final evaluation, only one training run was available for each configuration, and no independent simulated or experimental test set was used. Repeated training with controlled random seeds, a separate final test set, and validation on real trajectory measurements would be required for stronger claims about generalization and practical applicability. Within these limitations, the experiments show that GRUs can recover several latent physical parameters from trajectory data, provided that the training distribution reflects the expected degree of trajectory incompleteness and measurement noise.
